\section{Gráficos}

Nesta seção iremos descrever brevemente gráficos para realizar análise univariada de
variáveis de acordo com o seu tipo.
Para variáveis qualitativas, os gráficos são bsicamente construídos com base na tabela de
frequência de cada classe, mostrando a contagem ou a frequência relativa em cada classe.
Alguns exemplos incluem:

\begin{itemize}
\item Gráfico de Barras -- uma barra é dedicada a cada classe, e o tamanho da barra é
proporcional à contagem de resgistros na classe.
O código \ref{cod:barra} mostra a criação de tal gráfico com a base \texttt{cd\_brasil}
usando o pacote \texttt{ggplot2} da linguagem R.

\begin{lstlisting}[language=R, caption={Gráfico de Barras}, label={cod:barra}]
library(ggplot2)

cd_brasil |>
    ggplot(aes(x = regiao)) +
        geom_bar(stat = "count")
\end{lstlisting}

\item Gráfico de Setores -- construído com um círculo de raio arbitrário em que a
contagem ou frequência de registros é proporcional à área do setor circular que o
representa.
O código \ref{cod:setores} constrói o gráfico de setores correspondente à distribuição
de \texttt{regiao} na base \texttt{cd\_brasil}.

\begin{lstlisting}[language=R, caption={Gráfico de Setores}, label={cod:setores}]
tab_regiao <- table(cd_brasil$regiao)

pie(tab_regiao, labels = names(tab_regiao))
\end{lstlisting}

\item Gráfico de Pirulito -- mesmo princípio que o de barras, mas com aspecto mais
simples, em que o tamanho é proporcional a uma linha vertical com um ponto no eixo do
valor da contagem ou frequência.
O código \ref{cod:pirulito} mostra a construção desse gráfico com \texttt{ggplot2},
combinando \texttt{geom\_segment} para a linha e \texttt{geom\_point} para o ponto no
topo.

\begin{lstlisting}[language=R, caption={Gráfico de Pirulito}, label={cod:pirulito}]
library(ggplot2)

cd_brasil |>
    group_by(regiao) |>
    summarise(count = n()) |>
    ggplot(aes(x = regiao, y = count)) +
        geom_segment(aes(xend = regiao, y = 0, yend = count)) +
        geom_point(size = 3)
\end{lstlisting}
\end{itemize}

Para variáveis numéricas discretas, barras e pirulito também podem funcionar se houver
um número não muito grande de realizações distintas no conjunto de dados.
Para variáveis contínuas, a discretização em intervalos permite gerar vários tipos de
gráficos que ajudam a entender a distribuição de valores da variável, como

\begin{itemize}
\item Histogramas -- construído também com barras com tamanho proporcional densidade
-- de registros ou frequência -- em cada intervalo. Há várias fórmulas para determinar
o número de intervalos (\textit{bins}), como a de Sturges, dada por
$k = 1 + 3.3 \log_{10}(n)$, da raíz quadrada, com $k = \sqrt{n}$, de Scott, com
$k = \frac{3.5 \sigma}{\sqrt[3]{n}}$, onde $\sigma$ é o desvio-padrão, ou a de
Freedman-Diaconis, dada por $k = \frac{2 \cdot \mathrm{IQR}}{\sqrt[3]{n}}$.
O código \ref{cod:histograma} constrói o histograma da variável \texttt{temp} da base
\texttt{cd\_poluicao}.

\begin{lstlisting}[language=R, caption={Histograma}, label={cod:histograma}]
hist(cd_poluicao$temp,
     xlab = "Temperatura",
     ylab = "Frequencia",
     border = "white",
     col = "darkblue")
\end{lstlisting}

\item Ramo-e-folha -- os ramos são uma coluna vertical com os valores interiros
distintos assumidos pelas realizações.
As folhas são as partes fracionárias, empilhadas verticalmente ou horizontalmente em seus
ramos.
O código \ref{cod:ramo-folha} mostra a construção do diagrama para a mesma variável
\texttt{temp}.

\begin{lstlisting}[language=R, caption={Ramo-e-folha}, label={cod:ramo-folha}]
stem(cd_poluicao$temp, scale = .5)
\end{lstlisting}

\item Gráfico Dispersão Unidimensional -- os dados são mapeados a pontos na reta
correspondentes a seus valores.
Valores repetidos podem ser empilhados, de modo a dar uma melhor visualização à distribuição
dos dados.
O código \ref{cod:dispersao 1d} mostra a construção desse gráfico, empilhando os pontos
repetidos com a opção \texttt{method = "stack"}.

\begin{lstlisting}[language=R, caption={Dispersão Unidimensional}, label={cod:dispersao 1d}]
stripchart(cd_poluicao$temp,
           xlab = "Temperatura",
           method = "stack",
           offset = 2,
           at = 0,
           pch = 19,
           cex = 0.5,
           col = "darkblue")
\end{lstlisting}
\end{itemize}
